\input{../tex-inputs/latex-leseansicht-vorspann}

\section[Arthur und Olga Schnitzler an Gustav Schwarzkopf, 18. 7. 1909]{L04516 Arthur und Olga Schnitzler an Gustav Schwarzkopf, 18. 7. 1909}
\nopagebreak\mylabel{L04516v}
\rehead{ }\normalsize\beginnumbering\briefempfaengerindex{Schwarzkopf, Gustav@\textsc{Schwarzkopf, Gustav}!zzzSchnitzler, Olga@\emph{von Olga Schnitzler}!1909-07-182@{18. 7. 1909}|(be}\briefempfaengerindex{Schwarzkopf, Gustav@\textsc{Schwarzkopf, Gustav}!zzzSchnitzler, Arthur@\emph{von Arthur Schnitzler}!1909-07-182@{18. 7. 1909}|(be}
\toendnotes[C]{\smallbreak\pagebreak[2]}
\correspDesc{Versand  durch Arthur Schnitzler, Olga Schnitzler am 18. 7. 1909 in Edlach
\newline{}Übermittlung  am 1909-07-19 in Edlach
\newline{}Erhalt  durch Gustav Schwarzkopf im Zeitraum [19. 7. 1909 – 23. 7. 1909?] in Wien}\toendnotes[C]{\smallbreak}
\Standort{DLA, A:Schnitzler, HS.1985.1.1897.}
\physDesc{Bildpostkarte, 267 Zeichen
\newline{}Handschrift Arthur Schnitzler: Bleistift, deutsche Kurrent
\newline{}Handschrift Olga Schnitzler: Bleistift, lateinische Kurrent
\newline{}Versand: Stempel: »\nobreak{}\oindex{Edlach@\textbf{Edlach}|pwk}Edlach Reichenau in N. Oe., 1\textcolor{gray}{9}/7 09, 8–\textcolor{gray}{1}2 \textcolor{gray}{V}\nobreak{}«.  }\toendnotes[C]{\smallbreak}
\pstart
           \noindent{}\centering{}{\pb}\textcolor{gray}{\textbf{Rax}}\oindex{Raxalpe@\textbf{Raxalpe}|pw}\hspace*{2.5em}\textcolor{gray}{\textbf{Edlach bei Reichenau}}\oindex{Edlach@\textbf{Edlach}|pw}\pend
           \pstart{}{\pb}\textsc{Herrn Gustav Schwarzkopf}\pend{}\pstart{}\textsc{Wien I\oindex{Wien@\textbf{Wien}|pw}}\pend{}\pstart{}\textsc{Tiefer Graben 23}\oindex{Wien@\textbf{Wien}!I., Innere Stadt@\textbf{I., Innere Stadt}!Tiefer Graben 23@\textbf{Tiefer Graben 23}|pw}.\pend{}{\bigskip}\vspace{1em}
\pstart
           18/7. 09\pend
           \vspace{0.5em}
\pstart
           das Wetter wird beſſer und \textcolor{gray}{am}{ }22. d. komt \textsc{Liesl\pwindex{Steinrück, Elisabeth 19.\,11.\,1885 – 7.\,4.\,1920 Partenkirchen@\textsc{Steinrück, Elisabeth} (19.\,11.\,1885 – 7.\,4.\,1920 Partenkirchen)|pw}} (auf
               14 Tage.) der Bub\pwindex{Schnitzler, Heinrich 9.\,8.\,1902 Hinterbrühl – 12.\,7.\,1982 Wien@\textsc{Schnitzler, Heinrich} (9.\,8.\,1902 Hinterbrühl – 12.\,7.\,1982 Wien)|pwv} am
      Dienſtag. Und Sie?\pend
           \pstart Herzlichſt Ihr \spacefill\mbox{Arthur}\pend{}\selectlanguage{ngerman}\vspace{1em}
\pstart
           \noindent{}{[}hs. O. Schnitzler:{]} Ich war doch noch in Wien\oindex{Wien@\textbf{Wien}|pw} beim Heini\pwindex{Schnitzler, Heinrich 9.\,8.\,1902 Hinterbrühl – 12.\,7.\,1982 Wien@\textsc{Schnitzler, Heinrich} (9.\,8.\,1902 Hinterbrühl – 12.\,7.\,1982 Wien)|pw}.
      Kommen Sie{[},{]} die Lisl\pwindex{Steinrück, Elisabeth 19.\,11.\,1885 – 7.\,4.\,1920 Partenkirchen@\textsc{Steinrück, Elisabeth} (19.\,11.\,1885 – 7.\,4.\,1920 Partenkirchen)|pw} u andre Leute würden sich 
      sehr freuen!\pend
           \pstart \spacefill\mbox{Olga Schnitzler.}\pend{}\selectlanguage{ngerman}\endnumbering\briefempfaengerindex{Schwarzkopf, Gustav@\textsc{Schwarzkopf, Gustav}!zzzSchnitzler, Olga@\emph{von Olga Schnitzler}!1909-07-182@{18. 7. 1909}|)be}\briefempfaengerindex{Schwarzkopf, Gustav@\textsc{Schwarzkopf, Gustav}!zzzSchnitzler, Arthur@\emph{von Arthur Schnitzler}!1909-07-182@{18. 7. 1909}|)be}\mylabel{L04516h}
\begin{anhang}
\end{anhang}\newcommand{\dateiname}{L04516}\newcommand{\titel}{Arthur und Olga Schnitzler an Gustav Schwarzkopf, 18. 7. 1909}\newcommand{\editorInnen}{Herausgegeben von Jahnke, SelmaMüller, Martin Anton}\input{../tex-inputs/latex-leseansicht-abspann}
